\documentclass[11pt,a4paper]{article}

% --- Fonts & Encoding ---
\usepackage[T1]{fontenc}
\usepackage[utf8]{inputenc}
\usepackage{newtxtext,newtxmath}
\usepackage{microtype}

% --- Geometry & Spacing ---
\usepackage[margin=1in]{geometry}
\emergencystretch=1em
\raggedbottom

% ==== LINK COLOR TOGGLE ======================================
% Flip this switch:
\newif\ifusecaltechcolor
%\usecaltechcolortrue      % <-- Caltech Orange ON
 \usecaltechcolorfalse    % <-- Black links (no accent)

\usepackage{xcolor}
\definecolor{CaltechOrange}{HTML}{FF6C0C}
\definecolor{SoftGray}{HTML}{666666}

\usepackage[colorlinks=true]{hyperref}

\ifusecaltechcolor
  \hypersetup{linkcolor=CaltechOrange,urlcolor=CaltechOrange,citecolor=CaltechOrange}
  \newcommand{\RuleColor}{CaltechOrange}
\else
  \hypersetup{linkcolor=black,urlcolor=black,citecolor=black}
  \newcommand{\RuleColor}{black}
\fi
% =============================================================

% --- Discreet continuation header and page count ---
\usepackage{fancyhdr,lastpage}
\setlength{\headheight}{14pt}
\pagestyle{fancy}
\fancyhf{}
\fancyhead[L]{\footnotesize\color{SoftGray}Antonio Anna Mele}
\fancyhead[R]{\footnotesize\color{SoftGray}Publications}
\fancyfoot[C]{\footnotesize\color{SoftGray}\thepage\,/\,\pageref*{LastPage}}
\renewcommand{\headrulewidth}{0.3pt}
\renewcommand{\footrulewidth}{0pt}
\fancypagestyle{firstpage}{%
  \fancyhead{}%
  \renewcommand{\headrulewidth}{0pt}%
}

% --- Enumerated list with hanging indent ---
\usepackage{enumitem}
\setlist[enumerate]{%
  label=[\arabic*],leftmargin=2.3em,labelsep=0.6em,
  itemsep=10.8pt,topsep=5pt,parsep=0pt,partopsep=0pt,
  before={\fontsize{11.5}{15}\selectfont\interlinepenalty=10000}%
}
\newcommand{\cvpubitem}[1]{\item #1}
\newcommand{\authornote}[1]{%
  \unskip\ \mbox{\fontsize{9.5}{11.5}\selectfont\color{SoftGray}\textit{(#1)}}%
}

% --- Header (no duplicate section title) ---
\newcommand{\Header}{%
  {\fontsize{18}{21}\selectfont\bfseries All Publications (including preprints)}\par
  \vspace{2pt}
  {\normalsize \textit{Antonio Anna Mele --- Freie Universit\"at Berlin}}\par
  \vspace{2pt}
  {\footnotesize\color{SoftGray}%
  2025 Google PhD Fellow (Quantum Computing) \,|\,
  \href{mailto:antoniomele.p@gmail.com}{antoniomele.p@gmail.com} \,|\,
  \href{https://antonioannamele.com/research/}{Website} \,|\,
  \href{https://scholar.google.com/citations?user=IgnTVy4AAAAJ&hl=en}{Google Scholar}
  }\par
  \vspace{3pt}
  {\footnotesize\color{SoftGray}Listed in reverse chronological order by publication, acceptance, or initial preprint date.}\par
  \vspace{0.6em}\color{\RuleColor}\hrule height0.5pt\color{black}\vspace{0.4em}
}

% List sorted by most recent formal milestone, reverse chronologically.
% Journal articles: first official online publication date.
% Unpublished preprints: first arXiv posting date.
% Nature Physics (accepted, no published DOI): September 2026 acceptance, as documented
% on the author's news page. It is placed between late-September and August entries.
% No separation of preprints and published articles in the visible document.
\begin{document}
\thispagestyle{firstpage}
\Header

\begin{enumerate}

% ======= 2026 =======
% 2026-10-01: Adaptivity is all you need: Optimal stabilizer learning using just single-copy measurements
\cvpubitem{L.~Bittel, J.~Eisert, W.~Gong, \textbf{A.~A.~Mele}$^{\dagger}$, L.~Schatzki.
\emph{Adaptivity is all you need: Optimal stabilizer learning using just single-copy measurements}.
\href{https://arxiv.org/abs/2610.02031}{arXiv:2610.02031} (2026).
\authornote{$\dagger$ Authors listed alphabetically.}}

% 2026-09-30: Testing quantum Gaussianity with constant sample complexity
\cvpubitem{M.~Yaghubi Rad, R.~Puig, S.~Hassanandani, L.~Coffman, J.~Eisert, C.~Bravo-Prieto, \textbf{A.~A.~Mele}.
\emph{Testing quantum Gaussianity with constant sample complexity}.
\href{https://arxiv.org/abs/2609.40270}{arXiv:2609.40270} (2026).}

% 2026-09-29: Towards verifiable quantum advantage with random circuits: Observables that survive concentration
\cvpubitem{\textbf{A.~A.~Mele}$^{*}$, F.~A.~Mele$^{*}$, J.~R.~McClean, T.~E.~O'Brien.
\emph{Towards verifiable quantum advantage with random circuits: Observables that survive concentration}.
\href{https://arxiv.org/abs/2609.37890}{arXiv:2609.37890} (2026).
\authornote{$^{*}$ Equal contributors.}}

% 2026-09-22: When are bosonic Gaussian states classical to learn?
\cvpubitem{S.~Chen, \textbf{A.~A.~Mele}$^{\dagger}$, F.~A.~Mele, J.~Preskill.
\emph{When are bosonic Gaussian states classical to learn?}
\href{https://arxiv.org/abs/2609.26705}{arXiv:2609.26705} (2026).
\authornote{$\dagger$ Authors listed alphabetically.}}

% 2026-09-18: Optimal learning of quantum channels in diamond distance
\cvpubitem{\textbf{A.~A.~Mele}, L.~Bittel.
\emph{Optimal learning of quantum channels in diamond distance}.
\textit{Nature Physics} (accepted for publication, 2026). \href{https://arxiv.org/abs/2512.10214}{arXiv:2512.10214}.}

% 2026-08-13: Strong unitary designs in optimal depth and space
\cvpubitem{T.~Parella-Dilm\'e, J.~Barber\`a-Rodr\'iguez, S.~F.~E.~Oliviero, \textbf{A.~A.~Mele}.
\emph{Strong unitary designs in optimal depth and space}.
\href{https://arxiv.org/abs/2608.13491}{arXiv:2608.13491} (2026).}

% 2026-07-22: A complete theory of the Clifford commutant
\cvpubitem{L.~Bittel, J.~Eisert, L.~Leone, \textbf{A.~A.~Mele}$^{\dagger}$, S.~F.~E.~Oliviero.
\emph{A complete theory of the Clifford commutant}.
\href{https://quantum-journal.org/papers/q-2026-07-22-2171/}{\textit{Quantum} 10, 2171 (2026)}.
\authornote{$\dagger$ Authors listed alphabetically.}}

% 2026-07-15: Quantum memory advantage for quantum process tomography
\cvpubitem{C.~Bravo-Prieto, W.~Gong, \textbf{A.~A.~Mele}$^{\dagger}$.
\emph{Quantum memory advantage for quantum process tomography}.
\href{https://arxiv.org/abs/2607.13476}{arXiv:2607.13476} (2026).
\authornote{$\dagger$ Authors listed alphabetically.}}

% 2026-06-18: Experimental measurement and a physical interpretation of quantum shadow enumerators
\cvpubitem{D.~Miller, K.~Levi, L.~Postler, A.~Steiner, L.~Bittel, G.~A.~L.~White, Y.~Tang, E.~J.~Kuehnke, \textbf{A.~A.~Mele}, et~al.
\emph{Experimental measurement and a physical interpretation of quantum shadow enumerators}.
\href{https://doi.org/10.1103/8h3b-brg1}{\textit{Physical Review Research} 8, 023318 (2026)}. Editors' Suggestion.}

% 2026-04-22: Simulating quantum circuits with arbitrary local noise using Pauli propagation
\cvpubitem{A.~Angrisani, \textbf{A.~A.~Mele}, M.~S.~Rudolph, M.~Cerezo, Z.~Holmes.
\emph{Simulating quantum circuits with arbitrary local noise using Pauli propagation}.
\href{https://doi.org/10.1103/fb28-wlv2}{\textit{PRX Quantum} 7, 020313 (2026)}.}

% 2026-04-02: Noise-induced shallow circuits and the absence of barren plateaus
\cvpubitem{\textbf{A.~A.~Mele}, A.~Angrisani, S.~Ghosh, S.~Khatri, J.~Eisert, D.~S.~Fran\c{c}a, Y.~Quek.
\emph{Noise-induced shallow circuits and the absence of barren plateaus}.
\href{https://www.nature.com/articles/s41567-026-03245-z}{\textit{Nature Physics} 22, 751--756 (2026)}.}

% 2026-02-05: Analyzing the free states of one quantum resource theory as resource states of another
\cvpubitem{A.~E.~Deneris, P.~Braccia, P.~Bermejo, N.~L.~Diaz, \textbf{A.~A.~Mele}, M.~Cerezo.
\emph{Analyzing the free states of one quantum resource theory as resource states of another}.
\href{https://doi.org/10.1002/qute.202500702}{\textit{Advanced Quantum Technologies} 9, e00702 (2026)}.}


% ======= 2025 =======
% 2025-11-26: Learning quantum states of continuous-variable systems
\cvpubitem{F.~A.~Mele, \textbf{A.~A.~Mele}, L.~Bittel, J.~Eisert, V.~Giovannetti, L.~Lami, L.~Leone, S.~F.~E.~Oliviero.
\emph{Learning quantum states of continuous-variable systems}.
\href{https://www.nature.com/articles/s41567-025-03086-2}{\textit{Nature Physics} 21, 2002--2008 (2025)}.}

% 2025-10-07: Efficient learning of bosonic Gaussian unitaries
\cvpubitem{M.~Fanizza, V.~Iyer, J.~Lee, \textbf{A.~A.~Mele}$^{\dagger}$, F.~A.~Mele.
\emph{Efficient learning of bosonic Gaussian unitaries}.
\href{https://arxiv.org/abs/2510.05531}{arXiv:2510.05531} (2025).
\authornote{$\dagger$ Authors listed alphabetically.}}

% 2025-09-03: Optimal trace-distance bounds for free-fermionic states: Testing and improved tomography
\cvpubitem{L.~Bittel, \textbf{A.~A.~Mele}, J.~Eisert, L.~Leone.
\emph{Optimal trace-distance bounds for free-fermionic states: Testing and improved tomography}.
\href{https://doi.org/10.1103/pzx6-nkfb}{\textit{PRX Quantum} 6, 030341 (2025)}.}

% 2025-08-20: Energy-independent tomography of Gaussian states
\cvpubitem{L.~Bittel, F.~A.~Mele, J.~Eisert, \textbf{A.~A.~Mele}.
\emph{Energy-independent tomography of Gaussian states}.
\href{https://arxiv.org/abs/2508.14979}{arXiv:2508.14979} (2025).}

% 2025-07-14: A unified approach to quantum resource theories and a new class of free operations
\cvpubitem{N.~L.~Diaz, \textbf{A.~A.~Mele}, P.~Bermejo, P.~Braccia, A.~E.~Deneris, M.~Larocca, M.~Cerezo.
\emph{A unified approach to quantum resource theories and a new class of free operations}.
\href{https://arxiv.org/abs/2507.10851}{arXiv:2507.10851} (2025).}

% 2025-06-24: Characterizing quantum resourcefulness via group-Fourier decompositions
\cvpubitem{P.~Bermejo, P.~Braccia, \textbf{A.~A.~Mele}, N.~L.~Diaz, A.~E.~Deneris, M.~Larocca, M.~Cerezo.
\emph{Characterizing quantum resourcefulness via group-Fourier decompositions}.
\href{https://arxiv.org/abs/2506.19696}{arXiv:2506.19696} (2025).}

% 2025-06-12: Real classical shadows
\cvpubitem{M.~West, \textbf{A.~A.~Mele}, M.~Larocca, M.~Cerezo.
\emph{Real classical shadows}.
\href{https://doi.org/10.1088/1751-8121/addfc2}{\textit{Journal of Physics A: Mathematical and Theoretical} 58, 245304 (2025)}.}

% 2025-06-12: Optimal estimates of trace distance between bosonic Gaussian states and applications to learning
\cvpubitem{L.~Bittel, F.~A.~Mele, \textbf{A.~A.~Mele}, S.~Tirone, L.~Lami.
\emph{Optimal estimates of trace distance between bosonic Gaussian states and applications to learning}.
\href{https://quantum-journal.org/papers/q-2025-06-12-1769/}{\textit{Quantum} 9, 1769 (2025)}.}

% 2025-05-22: Stochastic noise can be helpful for variational quantum algorithms
\cvpubitem{J.~Liu, F.~Wilde, \textbf{A.~A.~Mele}, X.~Jin, L.~Jiang, J.~Eisert.
\emph{Stochastic noise can be helpful for variational quantum algorithms}.
\href{https://doi.org/10.1103/PhysRevA.111.052441}{\textit{Physical Review A} 111, 052441 (2025)}.}

% 2025-03-20: PAC-learning of free-fermionic states is NP-hard
\cvpubitem{L.~Bittel, \textbf{A.~A.~Mele}, J.~Eisert, L.~Leone.
\emph{PAC-learning of free-fermionic states is NP-hard}.
\href{https://quantum-journal.org/papers/q-2025-03-20-1665/}{\textit{Quantum} 9, 1665 (2025)}.}

% 2025-01-28: Efficient learning of quantum states prepared with few fermionic non-Gaussian gates
\cvpubitem{\textbf{A.~A.~Mele}, Y.~Herasymenko.
\emph{Efficient learning of quantum states prepared with few fermionic non-Gaussian gates}.
\href{https://doi.org/10.1103/PRXQuantum.6.010319}{\textit{PRX Quantum} 6, 010319 (2025)}.}


% ======= 2024 =======
% 2024-12-12: Learning fermionic correlations by evolving with random translationally invariant Hamiltonians
\cvpubitem{J.~Denzler, \textbf{A.~A.~Mele}, E.~Derbyshire, T.~Guaita, J.~Eisert.
\emph{Learning fermionic correlations by evolving with random translationally invariant Hamiltonians}.
\href{https://doi.org/10.1103/PhysRevLett.133.240604}{\textit{Physical Review Letters} 133, 240604 (2024)}.}

% 2024-09-24: Random ensembles of symplectic and unitary states are indistinguishable
\cvpubitem{M.~West, \textbf{A.~A.~Mele}, M.~Larocca, M.~Cerezo.
\emph{Random ensembles of symplectic and unitary states are indistinguishable}.
\href{https://arxiv.org/abs/2409.16500}{arXiv:2409.16500} (2024).}

% 2024-05-08: Introduction to Haar measure tools in quantum information: A beginner's tutorial
\cvpubitem{\textbf{A.~A.~Mele}.
\emph{Introduction to Haar measure tools in quantum information: A beginner's tutorial}.
\href{https://quantum-journal.org/papers/q-2024-05-08-1340/}{\textit{Quantum} 8, 1340 (2024)}.}


% ======= 2023 =======
% 2023-03-17: Exploiting symmetry in variational quantum machine learning
\cvpubitem{J.~J.~Meyer, M.~Mularski, E.~Gil-Fuster, \textbf{A.~A.~Mele}, F.~Arzani, A.~Wilms, J.~Eisert.
\emph{Exploiting symmetry in variational quantum machine learning}.
\href{https://doi.org/10.1103/PRXQuantum.4.010328}{\textit{PRX Quantum} 4, 010328 (2023)}.}


% ======= 2022 =======
% 2022-12-19: Avoiding barren plateaus via transferability of smooth solutions in Hamiltonian variational ansatz
\cvpubitem{\textbf{A.~A.~Mele}, G.~B.~Mbeng, G.~E.~Santoro, M.~Collura, P.~Torta.
\emph{Avoiding barren plateaus via transferability of smooth solutions in Hamiltonian variational ansatz}.
\href{https://doi.org/10.1103/PhysRevA.106.L060401}{\textit{Physical Review A} 106, L060401 (2022)}.}

\end{enumerate}
\end{document}
